\section{2017级工科数学分析下 A卷}

\subsection{试题}

\paragraph*{一、填空题：共 5 题，每题 2 分，共 10 分}
\begin{enumerate}
  \item 微分方程 $y''+y'-2y=1-2x$ 的通解为 \underline{\hspace{5cm}}。

  \item 设函数 $u=\ln(x^2+y^2+z^2)$，求 $\operatorname{div}(\operatorname{grad}u)=$ \underline{\hspace{4cm}}。

  \item 设 $\Gamma$ 是球面 $x^2+y^2+z^2=R^2$ 与平面 $x+y+z=0$ 的交线，则第一类曲线积分
  \[
    \oint_{\Gamma} y^2\,ds=
  \]
  \underline{\hspace{4cm}}。

  \item 参数曲线
  \[
    \begin{cases}
      x=t-\cos t,\\
      y=1-\sin t,\\
      z=t,
    \end{cases}
  \]
  在 $t=0$ 对应的点处的切线方程为 \underline{\hspace{6cm}}。

  \item 设周期为 $2\pi$ 的函数
  \[
    f(x)=
    \begin{cases}
      -1, & -\pi<x\le 0,\\
      1, & 0<x\le \pi,
    \end{cases}
  \]
  则 $f(x)$ 的傅里叶（Fourier）级数在 $x=\pi$ 处收敛于 \underline{\hspace{4cm}}。
\end{enumerate}

\paragraph*{二、单选题：共 5 题，每题 2 分，共 10 分}
\begin{enumerate}
  \item 关于未知函数 $y$ 的微分方程
  \[
    (y-\sin x)\,dx+\ln x\,dy=0
  \]
  是（\quad）。
  \[
  \begin{array}{ll}
    \text{A. 可分离变量方程；} & \text{B. 一阶线性非齐次方程；}\\
    \text{C. 一阶线性齐次方程；} & \text{D. 非线性方程。}
  \end{array}
  \]

  \item 二元函数
  \[
    f(x,y)=
    \begin{cases}
      x\sin\dfrac{1}{y}, & xy\ne 0,\\
      0, & xy=0,
    \end{cases}
  \]
  则 $\displaystyle \lim_{(x,y)\to(0,0)}f(x,y)$（\quad）。
  \[
  \begin{array}{llll}
    \text{A. 不存在；} & \text{B. 等于 1；} &
    \text{C. 等于 0；} & \text{D. 等于 2。}
  \end{array}
  \]

  \item 设函数 $z=f(x,y)$ 在 $(a,b)$ 的某个邻域内有直到二阶的连续偏导数，且
  \[
    \frac{\partial z}{\partial x}(a,b)=0,\qquad
    \frac{\partial z}{\partial y}(a,b)=0.
  \]
  记
  \[
    A=\frac{\partial^2 z}{\partial x^2}(a,b),\qquad
    B=\frac{\partial^2 z}{\partial x\partial y}(a,b),\qquad
    C=\frac{\partial^2 z}{\partial y^2}(a,b).
  \]
  则函数 $z=f(x,y)$ 在点 $(a,b)$ 取极大值的充分条件是（\quad）。
  \[
  \begin{array}{ll}
    \text{A. } A>0,\ AC>B^2; & \text{B. } A<0,\ AC>B^2;\\
    \text{C. } A>0,\ AC<B^2; & \text{D. } A<0,\ AC<B^2.
  \end{array}
  \]

  \item 函数 $\ln(1-x)$ 在 $x=0$ 处的泰勒（Taylor）展开式正确的是（\quad）。
  \[
  \begin{array}{ll}
    \text{A. } \displaystyle \ln(1-x)=\sum_{n=1}^{\infty}\frac{(-1)^n}{n}x^n,\quad x\in(-1,1];&
    \text{B. } \displaystyle \ln(1-x)=\sum_{n=1}^{\infty}\frac{(-1)^{n-1}}{n}x^n,\quad x\in(-1,1];\\[1.2em]
    \text{C. } \displaystyle \ln(1-x)=-\sum_{n=1}^{\infty}\frac{1}{n}x^n,\quad x\in[-1,1);&
    \text{D. } \displaystyle \ln(1-x)=\sum_{n=1}^{\infty}\frac{1}{n}x^n,\quad x\in[-1,1).
  \end{array}
  \]

  \item 使得级数
  \[
    \sum_{n=1}^{\infty}\frac{(-1)^n\ln n}{n^p}
  \]
  条件收敛的常数 $p$ 的取值范围是（\quad）。
  \[
  \begin{array}{llll}
    \text{A. }p\le 0; & \text{B. }0<p<1; &
    \text{C. }0<p\le 1; & \text{D. }p>1.
  \end{array}
  \]
\end{enumerate}

\paragraph*{三、计算题：共 3 题，每题 10 分，共 30 分}
\begin{enumerate}
  \item 设 $u=f(x,y,z)$，其中函数 $f$ 有二阶连续的偏导数，且 $z=z(x,y)$ 由方程
  \[
    z^5-5xy+5z=1
  \]
  所确定，求 $\dfrac{\partial u}{\partial x}$ 和 $\dfrac{\partial^2u}{\partial x^2}$。

  \item 计算累次积分
  \[
    \int_{\frac14}^{\frac12} dx\int_{\frac12}^{\sqrt{x}} e^{x/y}\,dy
    +\int_{\frac12}^{1} dx\int_x^{\sqrt{x}} e^{x/y}\,dy.
  \]

  \item 计算球面 $x^2+y^2+z^2=4z$ 含在抛物面 $z=x^2+y^2$ 内的部分的面积。
\end{enumerate}

\paragraph*{四、解答题：共 3 题，每题 10 分，共 30 分}
\begin{enumerate}
  \item 设 $\Sigma$ 是锥面 $z=\sqrt{x^2+y^2}$ 被平面 $z=0$ 及 $z=1$ 截下的部分的下侧，计算第二类曲面积分
  \[
    \iint_{\Sigma} xye^z\,dy\,dz+yz^2\,dz\,dx-ye^z\,dx\,dy.
  \]

  \item 设曲线积分
  \[
    \int_{\Gamma}\left(\sin x-f(x)\right)\frac{y}{x}\,dx+f(x)\,dy
  \]
  与路径无关，其中 $f(x)$ 有一阶连续导数且 $f(\pi)=0$，求 $f(x)$ 并计算曲线积分
  \[
    \int_{(1,0)}^{(\pi,\pi)}
    \left(\sin x-f(x)\right)\frac{y}{x}\,dx+f(x)\,dy.
  \]

  \item 求幂级数
  \[
    \sum_{n=0}^{\infty}\frac{1}{2n+1}x^{2n+1}
  \]
  的收敛域及和函数。
\end{enumerate}

\paragraph*{五、证明题：共 1 题，每题 10 分，共 10 分}
证明函数项级数
\[
  \sum_{n=0}^{\infty}x^2e^{-nx}
\]
在 $[0,+\infty)$ 一致收敛。

\paragraph*{六、应用题：共 1 题，每题 10 分，共 10 分}
将长度为 $2a$ 的铁丝分成三段，分别围成一个正方形、一个圆形和一个正三角形，求三个图形面积之和的最小值。

\subsection{答案与解析}

\paragraph*{一、填空题}
\begin{enumerate}
  \item
  \[
    y=C_1e^{-2x}+C_2e^x+x.
  \]

  \item
  \[
    \operatorname{div}(\operatorname{grad}u)
    =\Delta u=\frac{2}{x^2+y^2+z^2}.
  \]

  \item
  \[
    \oint_{\Gamma}y^2\,ds=\frac{2}{3}\pi R^3.
  \]

  \item 当 $t=0$ 时，点为 $(-1,1,0)$，切向量为 $(1,-1,1)$，故切线方程为
  \[
    \frac{x+1}{1}=\frac{y-1}{-1}=\frac{z}{1}.
  \]

  \item 傅里叶级数在跳跃点收敛于左右极限的平均值。由周期延拓知
  \[
    \frac{f(\pi-0)+f(\pi+0)}{2}=\frac{1+(-1)}{2}=0.
  \]
\end{enumerate}

\paragraph*{二、单选题}
\begin{enumerate}
  \item 方程可化为 \(y'+\dfrac1{\ln x}y=\dfrac{\sin x}{\ln x}\)，是一阶线性非齐次方程。选 B。
  \item 因 \(|x\sin(1/y)|\le |x|\)，故 $(x,y)\to(0,0)$ 时极限为 $0$。选 C。
  \item 二元函数极大值的充分条件为 \(A<0,\ AC-B^2>0\)。选 B。
  \item 按题面函数应有 \(\displaystyle \ln(1-x)=-\sum_{n=1}^{\infty}\frac{x^n}{n}\)，收敛区间为 $[-1,1)$，故选 C。\textit{（校勘：答案页标为 A，与题面函数及端点范围不符。）}
  \item \(\displaystyle \sum (-1)^n\frac{\ln n}{n^p}\) 在 $p>0$ 时由交错判别法收敛，在 $p>1$ 时绝对收敛，故条件收敛范围为 $0<p\le1$。选 C。
\end{enumerate}

\paragraph*{三、计算题}
\begin{enumerate}
  \item 对方程 $z^5-5xy+5z=1$ 两边求全微分，得
  \[
    5z^4\,dz-5x\,dy-5y\,dx+5\,dz=0.
  \]
  因此
  \[
    dz=\frac{y\,dx+x\,dy}{z^4+1},
  \]
  即
  \[
    \frac{\partial z}{\partial x}=\frac{y}{z^4+1},
    \qquad
    \frac{\partial z}{\partial y}=\frac{x}{z^4+1}.
  \]
  由复合函数求导法则，
  \[
    \frac{\partial u}{\partial x}
    =f'_1(x,y,z)+f'_3(x,y,z)\frac{\partial z}{\partial x}
    =f'_1(x,y,z)+f'_3(x,y,z)\frac{y}{z^4+1}.
  \]
  进一步，
  \[
    \frac{\partial^2u}{\partial x^2}
    =f''_{11}(x,y,z)+2f''_{13}(x,y,z)\frac{y}{z^4+1}
    +f''_{33}(x,y,z)\left(\frac{y}{z^4+1}\right)^2
    -f'_3(x,y,z)\frac{4y^2z^3}{(z^4+1)^3}.
  \]

  \item 记 $D$ 为由 $y=\dfrac12$、$y=\sqrt{x}$ 和 $y=x$ 围成的区域。原积分化为
  \[
    \iint_D e^{x/y}\,dx\,dy
    =\int_{\frac12}^{1}dy\int_{y^2}^{y}e^{x/y}\,dx.
  \]
  于是
  \[
    \int_{\frac12}^{1}dy\int_{y^2}^{y}e^{x/y}\,dx
    =\int_{\frac12}^{1}\left(ey-ye^y\right)\,dy
    =\frac{3}{8}e-\frac12e^{1/2}.
  \]

  \item 球面可写为
  \[
    z=2+\sqrt{4-x^2-y^2}.
  \]
  其含在抛物面 $z=x^2+y^2$ 内的部分在 $xOy$ 平面上的投影区域为
  \[
    D=\{(x,y)\mid x^2+y^2\le 3\}.
  \]
  曲面面积微元为
  \[
    dS=\sqrt{1+\left(\frac{\partial z}{\partial x}\right)^2
    +\left(\frac{\partial z}{\partial y}\right)^2}\,dx\,dy
    =\frac{2}{\sqrt{4-x^2-y^2}}\,dx\,dy.
  \]
  因此所求面积为
  \[
    A=\iint_D\frac{2}{\sqrt{4-x^2-y^2}}\,dx\,dy
    =\int_0^{2\pi}d\theta\int_0^{\sqrt3}\frac{2r}{\sqrt{4-r^2}}\,dr
    =4\pi.
  \]
\end{enumerate}

\paragraph*{四、解答题}
\begin{enumerate}
  \item 考虑平面 $\Sigma_1:z=1$ 在 $x^2+y^2\le 1$ 的部分，取上侧，则 $\Sigma\cup\Sigma_1$ 形成闭曲面，取外侧。记其围成的闭区域为 $\Omega$。由 Gauss 公式，
  \[
    \iint_{\Sigma\cup\Sigma_1} xye^z\,dy\,dz+yz^2\,dz\,dx-ye^z\,dx\,dy
    =\iiint_{\Omega}z^2\,dx\,dy\,dz.
  \]
  而
  \[
    \iiint_{\Omega}z^2\,dx\,dy\,dz
    =\int_0^1 dz\iint_{x^2+y^2\le z^2}z^2\,dx\,dy
    =\int_0^1\pi z^4\,dz=\frac{\pi}{5}.
  \]
  又 $\Sigma_1$ 的单位正法向量为 $\mathbf n=(0,0,1)$，故
  \[
    \iint_{\Sigma_1} xye^z\,dy\,dz+yz^2\,dz\,dx-ye^z\,dx\,dy
    =\iint_{x^2+y^2\le 1}(-ye)\,dx\,dy=0.
  \]
  因此
  \[
    \iint_{\Sigma} xye^z\,dy\,dz+yz^2\,dz\,dx-ye^z\,dx\,dy
    =\frac{\pi}{5}.
  \]

  \item 记
  \[
    P=\left(\sin x-f(x)\right)\frac{y}{x},\qquad Q=f(x).
  \]
  由曲线积分与路径无关，有 $\dfrac{\partial P}{\partial y}=\dfrac{\partial Q}{\partial x}$，即
  \[
    \frac{\sin x}{x}-\frac{f(x)}{x}=f'(x).
  \]
  这是一个一阶线性非齐次微分方程，积分因子为 $x$，从而
  \[
    (xf(x))'=xf'(x)+f(x)=\sin x.
  \]
  其通解为
  \[
    f(x)=-\frac{\cos x}{x}+\frac{C}{x}.
  \]
  由 $f(\pi)=0$ 得 $C=-1$。因此，
  \[
    f(x)=\frac{-\cos x-1}{x}.
  \]
  此时
  \[
    P\,dx+Q\,dy
    =d\left(\frac{(-\cos x-1)y}{x}\right).
  \]
  因此
  \[
    \int_{(1,0)}^{(\pi,\pi)}P\,dx+Q\,dy=0.
  \]

  \item 考虑幂级数
  \[
    \sum_{n=0}^{\infty}\frac{1}{2n+1}t^n.
  \]
  由于
  \[
    \lim_{n\to\infty}
    \frac{\frac{1}{2(n+1)+1}}{\frac{1}{2n+1}}
    =\lim_{n\to\infty}\frac{2n+1}{2n+3}=1,
  \]
  其收敛半径为 $1$。故原幂级数
  \[
    \sum_{n=0}^{\infty}\frac{1}{2n+1}x^{2n+1}
    =x\sum_{n=0}^{\infty}\frac{1}{2n+1}(x^2)^n
  \]
  的收敛半径也为 $1$。当 $x=\pm1$ 时，级数为 $\pm\sum_{n=0}^{\infty}\dfrac{1}{2n+1}$，发散。因此收敛域为 $(-1,1)$。

  设其和函数为
  \[
    f(x)=\sum_{n=0}^{\infty}\frac{1}{2n+1}x^{2n+1},\qquad x\in(-1,1).
  \]
  则
  \[
    f'(x)=\sum_{n=0}^{\infty}x^{2n}=\frac{1}{1-x^2}.
  \]
  因为 $f(0)=0$，所以
  \[
    f(x)=\int_0^x\frac{1}{1-t^2}\,dt
    =\frac12\ln\frac{1+x}{1-x},\qquad x\in(-1,1).
  \]
\end{enumerate}

\paragraph*{五、证明题}
证明一：对任意 $x\in[0,+\infty)$，当 $n\ge 1$ 时，
\[
  e^{nx}>1+nx+\frac12(nx)^2>\frac12(nx)^2.
\]
因此
\[
  |x^2e^{-nx}|=\frac{x^2}{e^{nx}}\le \frac{2}{n^2}.
\]
而数项级数 $\displaystyle \sum_{n=1}^{\infty}\frac{1}{n^2}$ 收敛。由优级数判别法，
\[
  \sum_{n=0}^{\infty}x^2e^{-nx}
\]
在 $[0,+\infty)$ 一致收敛。

证明二：该函数项级数的部分和为
\[
  S_n(x)=
  \begin{cases}
    \dfrac{x^2(1-e^{-(n+1)x})}{1-e^{-x}}, & x>0,\\
    0, & x=0.
  \end{cases}
\]
因此
\[
  S(x)=\lim_{n\to\infty}S_n(x)=
  \begin{cases}
    \dfrac{x^2}{1-e^{-x}}, & x>0,\\
    0, & x=0.
  \end{cases}
\]
下证 $S_n(x)$ 一致收敛于 $S(x)$。对任意 $\varepsilon>0$，取 $N>\dfrac1\varepsilon$，当 $n>N$ 时，对任意 $x\in[0,+\infty)$，
\[
  |S_n(x)-S(x)|
  \le \left|\frac{x^2e^{-(n+1)x}}{1-e^{-x}}\right|
  =\left|\frac{x}{e^x-1}\right|\left|\frac{x}{e^{nx}}\right|
  <\frac1n<\varepsilon.
\]
其中用到不等式 $e^x\ge x+1\ge x$。故该函数项级数在 $[0,+\infty)$ 一致收敛。

\paragraph*{六、应用题}
设正方形的边长为 $x$，圆形的半径为 $y$，正三角形的边长为 $z$，则三个图形的周长之和为
\[
  4x+2\pi y+3z=2a.
\]
它们的面积之和为
\[
  f(x,y,z)=x^2+\pi y^2+\frac{\sqrt3}{4}z^2.
\]
构造 Lagrange 函数
\[
  L(x,y,z,\lambda)
  =x^2+\pi y^2+\frac{\sqrt3}{4}z^2
  -\lambda(4x+2\pi y+3z-2a).
\]
极值点满足
\[
  \begin{cases}
    \dfrac{\partial L}{\partial x}=2x-4\lambda=0,\\
    \dfrac{\partial L}{\partial y}=2\pi y-2\pi\lambda=0,\\
    \dfrac{\partial L}{\partial z}=\dfrac{\sqrt3}{2}z-3\lambda=0,\\
    \dfrac{\partial L}{\partial \lambda}=-(4x+2\pi y+3z-2a)=0.
  \end{cases}
  \]
  解得
  \[
    \lambda=\frac{a}{4+\pi+3\sqrt3},\qquad
    x_0=\frac{2a}{4+\pi+3\sqrt3},\qquad
    y_0=\frac{a}{4+\pi+3\sqrt3},\qquad
    z_0=\frac{2\sqrt3a}{4+\pi+3\sqrt3}.
  \]
  所求最小面积为
  \[
    f(x_0,y_0,z_0)=\frac{a^2}{4+\pi+3\sqrt3}.
  \]
